\section{The stopping dichotomy and a conditional closure gate}
\label{sec:stopping}

We assemble the two quadraticizations.  For each positive divisor
$s\mid h$, put
\begin{equation}
 \cT_s^{\mathbb P}
 =\sum_{\substack{k\\(k,h)=s}}\beta_I(k)S_k.
 \label{eq:sector-tail}
\end{equation}
Write $\cC_{D,\mathrm{prim}}^{\mathrm{gen}}$ for the generic subsum
$\cC_D^{\mathrm{gen}}$ of \cref{eq:d-off,eq:generic-set} with the
fixed mask $\xi(d,j)=\one_{(dj,h)=1}$.

\begin{theorem}[Primitive/nonprimitive stopping dichotomy]
\label{thm:stopping-dichotomy}
Assume \eqref{eq:block-geometry}, $D_0<V/2$, and
$L\ge X^{\lambda_0}$ for a fixed $\lambda_0>0$.  Suppose that, for
fixed $a\ge0$ and $\eta>0$,
\begin{equation}
 |\cT_{L,K}(I)|\ge\frac{\eta X}{(\log X)^a}.
 \label{eq:tail-failure}
\end{equation}
Then, for every fixed $\kappa>0$ and all sufficiently large $X$, at
least one of the following alternatives holds.

\begin{enumerate}[label=\textup{(\Alph*)}]
\item \textbf{Primitive generic determinant witness.}
There is a dyadic $D\le V$, with $Q=LD$, such that the primitive
$d$-slice has a generic correlation satisfying
\begin{equation}
 \cC_{D,\mathrm{prim}}^{\mathrm{gen}}
 \gg_{\eta,h,W}
 \frac{XQ}{(\log X)^{2a+2}},
 \label{eq:primitive-generic-dichotomy}
\end{equation}
supported on
\[
 \ell_1\ne\ell_2,\qquad
 |\ell_1d_1-\ell_2d_2|>QX^{-\kappa},\qquad
 (d_1,d_2)\le X^\kappa,\qquad
 (d_1j,h)=(d_2j,h)=1.
\]

\item \textbf{Nonprimitive endpoint finite-model witness.}
There is a divisor $s>1$ of $h$ for which
\begin{equation}
 \cC_{I,s}^{MM,\mathrm{off}}
 \gg_{\eta,h,W}\frac{XL}{(\log X)^{2a}}.
 \label{eq:nonprimitive-dichotomy}
\end{equation}
Its M\"obius coefficients are supported only on the two endpoint
bands in \eqref{eq:endpoint-support}.  After a dyadic split of
$|\ell_1-\ell_2|$, one signed separation cell contributes
\begin{equation}
 \gg_{\eta,h,W}\frac{XL}{(\log X)^{2a+1}}.
 \label{eq:nonprimitive-Delta-cell}
\end{equation}
\end{enumerate}
\end{theorem}

\begin{proof}
By \cref{lem:source-primes}, the source prime-power error is smaller
than \eqref{eq:tail-failure} by a fixed power.  Since $(k,h)$ can take
only $O_h(1)$ values, one sector in \eqref{eq:sector-tail} has magnitude
$\gg_{\eta,h}X(\log X)^{-a}$.

If the large sector is $s=1$, open $\beta_I(k)$ and use a smooth dyadic
partition on $d\in I$.  There are $O(\log X)$ slices, so one primitive
slice \eqref{eq:TD}, with $\xi(d,j)=\one_{(dj,h)=1}$, has magnitude
$\gg X(\log X)^{-a-1}$.  Apply
\cref{thm:generic-stopping} with $a$ replaced by $a+1$ to obtain
\eqref{eq:primitive-generic-dichotomy}.

If the large sector has $s>1$, repeat the same-$k$ Cauchy argument of
\cref{thm:same-k-ttstar} with the restriction $(k,h)=s$.  Its diagonal
is still $O_\eps(X^{1+\eps})$, so
\[
 \cC_{I,s}^{\mathrm{off}}
 \gg_{\eta,h,W}XL(\log X)^{-2a}.
\]
The error in \eqref{eq:nonprimitive-MM} is negligible at this scale,
which proves \eqref{eq:nonprimitive-dichotomy}.  Endpoint localization
follows by squaring \eqref{eq:endpoint-collapse}.  There are
$O(\log X)$ nonempty dyadic source-separation ranges; since their signed
sum is positive and has the displayed size, at least one proves
\eqref{eq:nonprimitive-Delta-cell}.
\end{proof}

\begin{remark}[Block scope]
\Cref{thm:stopping-dichotomy} is attached to a symmetric tail block
for which the TPC-17 prefix theorem applies.  It does not assert that
every possible failure of a fixed-shift prime-pair asymptotic must occur
in this particular family of blocks.  That global statement requires
the full preceding normal-form interface.
\end{remark}

We now isolate the exact analytic input that would remove alternative
(A).  This input is deliberately stated rather than assumed true.

\begin{assumption}[Generic determinant dispersion
$\mathrm{DD}_h(\theta)$]
\label{input:DD}
Fix $\theta>1/2$.  For every fixed $\eps_0>0$, uniformly for every
relevant slice \eqref{eq:TD} with
$Q=LD\le X^{\theta-\eps_0}$, every fixed $\kappa>0$, every admissible
fixed-residue mask $\xi$ described in \cref{sec:opened-d}, and every
$A>0$, one has
\begin{equation}
 |\cC_D^{\mathrm{gen}}|
 \ll_{A,\eps_0,\kappa,h,W,\psi}
 XQ(\log X)^{-A}.
 \label{eq:DD}
\end{equation}
\end{assumption}

\begin{proposition}[Conditional slice and tail closure]
\label{prop:conditional-closure}
Assume the block geometry \eqref{eq:block-geometry}, so in particular
$L\ge X^{\lambda_0}$ for some fixed $\lambda_0>0$, and assume
$\mathrm{DD}_h(\theta)$ covers every member of the fixed bounded-overlap
dyadic partition of a given symmetric tail block.  Then, for every
$B>0$,
\begin{equation}
 \cT_{L,K}(I)\ll_BX(\log X)^{-B}.
 \label{eq:conditional-tail-closure}
\end{equation}
\end{proposition}

\begin{proof}
Combine \cref{thm:degenerate-removal,input:DD} with
\eqref{eq:d-first-TTstar}.  For any requested $B$, take the logarithmic
saving in \eqref{eq:DD} sufficiently large.  Since $J\asymp X/Q$,
\begin{align*}
 |\cT_D|
 &\ll \sqrt{JX^{1+\eps}}
     +\sqrt{JXQ}(\log X)^{-A/2}\\
 &\ll XQ^{-1/2}X^{\eps/2}+X(\log X)^{-A/2}.
\end{align*}
The first term has a fixed power saving because $Q\ge L\ge
X^{\lambda_0}$.  Sum the $O(\log X)$ divisor slices and restore the
negligible source prime powers using \cref{lem:source-primes}.
\end{proof}

\subsection{How strong is the conditional input?}

Write
\[
 L=X^{\lambda+o(1)},\qquad
 V=X^{v+o(1)},\qquad
 v=\frac14-\frac\delta2.
\]
The largest opened row has $Q=LV=X^{\lambda+v+o(1)}$.  Define the
boundary exponent
\begin{equation}
 \theta_{\mathrm{tail}}=\lambda+\frac14-\frac\delta2.
 \label{eq:theta-tail}
\end{equation}
Because \cref{input:DD} retains a fixed $\eps_0$ margin, full-tail
closure requires $\theta>\theta_{\mathrm{tail}}$ by a fixed amount;
equality is only the boundary value.  For the two explicit TPC-17
profiles, let $\sigma>0$ denote the fixed power-saving slack, subject
to $\sigma<1/100$ in the published row and $\sigma<1/1000$ in the
version-locked Li row.  The boundary ledger is:
\begin{center}
\begin{tabular}{@{}lll@{}}
\toprule
profile & $\lambda$ and allowed $\delta$ & range of $\theta_{\rm tail}$ \\
\midrule
published Maynard layer
& $10/21-\sigma$, $1/42+3\sigma<\delta<1/4$
& $101/168-\sigma$ to $5/7-(5/2)\sigma$ \\
Li v6 extension
& $8/17-\sigma$, $1/34+2\sigma<\delta<1/4$
& $81/136-\sigma$ to $12/17-2\sigma$ \\
\bottomrule
\end{tabular}
\end{center}
The endpoints in the table are limiting open endpoints and are not
attained by the displayed strict $\delta$ ranges.
Numerically the lower endpoints are about $0.6012$ and $0.5956$ when
$\sigma$ is suppressed, while the upper endpoints are about $0.7143$
and $0.7059$.  The Li row is version-locked to the unrefereed preprint
used by TPC-17; the published theorem chain uses only the first row.

Maynard's published well-factorable benchmark reaches level $3/5$ for
its specified coefficient class \citep{Maynard2025II}.  Lichtman's
preprint reaches $66/107\approx0.61682$ for triply well-factorable
weights, with a conditional $5/8$ statement under an additional
spectral conjecture \citep{Lichtman2023}.  The numerical overlap with
the lower end of \eqref{eq:theta-tail} does not imply applicability:
the sharp tail fails the balanced factor shape noted after
\cref{prop:weighted-second-moment}, and \eqref{eq:DD} is a two-residual
variable-determinant correlation rather than a one-prime progression
estimate.  Exponent size is not the missing hypothesis.

For clarity, the prime-pair singular series below is
\begin{equation}
 \mathfrak S(h)
 =\prod_p\frac{1-\nu_h(p)/p}{(1-1/p)^2},
 \qquad
 \nu_h(p)=\#\{a\bmod p:a(a+h)\equiv0\pmod p\}.
 \label{eq:singular-series}
\end{equation}
In particular it vanishes for odd $h$ and is positive for nonzero even
$h$.

\begin{corollary}[Why uniform closure reaches the fixed-shift problem]
\label{cor:DD-HL-boundary}
Assume the symmetric normal form of TPC-16 \citep{WangTPC16} and, as an additional global
cover hypothesis, that its hard packet admits
\begin{equation}
 \cH^{\mathrm{HB}}_{h,R;V}(X)
 =\sum_{b\in\mathcal B_X}(\cC_b+\cT_b)+o(X),
 \qquad |\mathcal B_X|=O_W(\log X),
 \label{eq:all-block-cover}
\end{equation}
where every block satisfies \eqref{eq:block-geometry}, the TPC-17
assembled-prefix estimate for $\cC_b$, and the opened-$d$ geometry for
$\cT_b$.  Assume also $\mathrm{DD}_h(\theta)$ with a fixed strict
margin above every boundary in \eqref{eq:theta-tail}.  Then the
resulting chain
gives the weighted fixed-shift Hardy--Littlewood asymptotic
\begin{equation}
 \sum_n\Lambda(n)\Lambda(n+h)W(n/X)
 =\mathfrak S(h)X I_0(W)+o(X).
 \label{eq:conditional-HL}
\end{equation}
For $h=2$ and nonnegative nonzero $W$, this would imply infinitely many
twin primes after the standard removal of prime powers.
\end{corollary}

\begin{proof}
The assumed all-block prefix estimates remove the $\cC_b$ terms.
Proposition \ref{prop:conditional-closure} removes each complementary
$\cT_b$.  The global cover \eqref{eq:all-block-cover} and the TPC-16
normal form identify the sum of those blocks with the difference
between the left side of \eqref{eq:conditional-HL} and its calibrated
singular-series main term.  There are only $O_W(\log X)$ blocks, so
arbitrary logarithmic saving survives their sum.  The $h=2$ consequence
follows from positivity and the fact that source or target prime powers
contribute $o(X)$.
\end{proof}

This corollary is a logical strength test, not a proof of its
hypotheses.  It explains why the primitive generic channel should not
be dismissed as a routine Kloosterman estimate.  By contrast,
\citet{MatomakiRadziwillTao2019} establish the expected prime
correlations for almost all shifts in sufficiently long ranges.  That
is a rigorous averaged-$h$ benchmark, but it supplies no direct
fixed-$h$, lcm-row-resolved estimate for \eqref{eq:DD}.
